\scanpage{106}
\deleted{inversible $\underline L$ sur $X$} Je dis qu'elle est partout définie. Soit $U$ son ouvert de définition, $T=\overline{\unsure{s}(U)}$, il suffit de prouver que $T\longrightarrow S$ [qui est \inserted{birationnel, et} surjectif d'après la propreté essentielle de $\underline{\operatorname{Pic}}_{X/S}$] est un isom, ou encore [\uuline{Main Theorem}] à fibres discrètes -- \unsure{ou si on veut}, fibres \inserted{géom} dénombrables sur $\uline{\forall}$ corps extension\ldots. Cela va résulter du fait suivant.

\uline{Lemme} Soit $t\in T$ au dessus de $s\in S$, \deleted{Alors} Soit $k$ une \deleted{clôture} extension algébriquement close de $k(t)$, et $\underline L_t$ le Module sur $X_k$ défini par $t$. Alors $L_t$ est définissable par un diviseur de Cartier de support $\subset Z_k$.

En eff, soit $V$ un anneau de val. discrète dominant $\mathcal O_{T,t}$, $S'=$
\margin{105}
